% !TeX root = ../Thesis.tex

%************************************************
\chapter{Ergänzende Informationen}\label{ch:ergänzendeinfos}
%************************************************
\glsresetall % Resets all acronyms to not used

% ============================================================
% Stempelgruppen
% ============================================================

\begin{landscape}

\small

\begin{longtable}{
    |c|
    |c|
    |c|
    |c|
    |p{2.5cm}|
    |p{2.2cm}|
    |p{2.5cm}|
    |p{2.5cm}|
}
    \caption{Übersicht der verwendeten Stempelgruppen}
    \label{tab:stempelgruppen} \\

    \hline
    \textbf{Nr.} &
    \textbf{Vorderseiten\-stempel} &
    \textbf{Münzen mit Vorderseiten\-bild} &
    \textbf{Verschiedene Rückseiten\-stempel} &
    \textbf{Münzstätte} &
    \textbf{Region} &
    \textbf{Datierung} &
    \textbf{Bild-/Objektart} \\
    \hline
    \endfirsthead

    \hline
    \textbf{Nr.} &
    \textbf{Vorderseiten\-stempel} &
    \textbf{Münzen mit Vorderseiten\-bild} &
    \textbf{Verschiedene Rückseiten\-stempel} &
    \textbf{Münzstätte} &
    \textbf{Region} &
    \textbf{Datierung} &
    \textbf{Bild-/Objektart} \\
    \hline
    \endhead

    1  & 973  & 9  & 5              & & & & \\ \hline
    2  & 985  & 9  & 8              & & & & \\ \hline
    3  & 1002 & 9  & 8              & & & & \\ \hline
    4  & 1249 & 10 & 6              & & & & \\ \hline
    5  & 1311 & 10 & 6              & & & & \\ \hline
    6  & 1386 & 10 & 6 (+1 ohne ID) & & & & \\ \hline
    7  & 463  & 10 & 6              & & & & \\ \hline
    8  & 1275 & 10 & 9              & & & & \\ \hline
    9  & 2785 & 10 & 6              & & & & \\ \hline
    10 & 2816 & 10 & 2              & & & & \\ \hline
    11 & 2905 & 10 & 6              & & & & \\ \hline
    12 & 2926 & 10 & 4 (+1 ohne ID) & & & & \\ \hline
    13 & 3570 & 7  & 5              & & & & \\ \hline
    14 & 3765 & 10 & 8              & & & & \\ \hline
    15 & 3788 & 9  & 8              & & & & \\ \hline
    16 & 3822 & 9  & 4              & & & & \\ \hline
    17 & 3829 & 10 & 5              & & & & \\ \hline
    18 & 960  & 10 & 7              & & & & \\ \hline
    19 & 1055 & 10 & 9              & & & & \\ \hline
    20 & 932  & 10 & 4              & & & & \\ \hline

\end{longtable}

\normalsize

\end{landscape}

% ============================================================
% Konfigurationsparameter
% ============================================================
Konfigurationsparameter:
\begin{table}[htbp]
    \centering
    \caption{Konfigurationsparameter der Anwendung}
    \label{tab:konfigurationsparameter}

    \begin{tabularx}{\textwidth}{|X|l|l|}
        \hline
        \textbf{Parameter} &
        \textbf{Wert} &
        \textbf{Datei} \\
        \hline

        Modell (verwendet)
            & \texttt{yolov8s-worldv2.pt}
            & \texttt{yolo\_config.json} \\
        \hline

        Modell (erster Ansatz)
            & \texttt{yolov8n.pt}
            & \texttt{yolo\_config.json} \\
        \hline

        Confidence-Schwellenwert
            & 0{,}005
            & \texttt{yolo\_config.json} \\
        \hline

        Anzahl der Prompts
            & 57 (Anhang B)
            & \texttt{yolo\_config.json} \\
        \hline

        Batchgröße (Ausschnitte)
            & 16
            & \texttt{data.json} \\
        \hline

        Anzahl Stempelgruppen
            & 20
            & \texttt{data.json} \\
        \hline

        Münzen je Stempelgruppe
            & 10
            & \texttt{data.json} \\
        \hline

        Seitengröße beim API-Abruf
            & 100
            & \texttt{data.json} \\
        \hline

        Positionstoleranz (Mittelpunktabstand)
            & 0{,}15
            & Objektkonfiguration \\
        \hline

        Schwellenwert Flächenverhältnis
            & 0{,}5
            & Objektkonfiguration \\
        \hline
    \end{tabularx}
\end{table}

\todo{TODO: Ultralytics-Standardwerte prüfen und anschließend ergänzen: Bildgröße: vermutlich 640 px, IoU-Schwelle NMS: vermutlich 0,7, Maximale Detektionen je Bild: vermutlich 300, verwendete Ultralytics-Version}

% ============================================================
% Im Code festgelegte Parameter
% ============================================================
Im Code festgelegte Parameter:
\begin{table}[htbp]
    \centering
    \caption{Im Programmcode festgelegte Parameter}
    \label{tab:codeparameter}

    \begin{tabularx}{\textwidth}{|X|X|}
        \hline
        \textbf{Parameter} &
        \textbf{Wert} \\
        \hline

        Schicht für die Feature-Vektoren
            & 21 \\
        \hline

        Standard-Ähnlichkeitsmaß
            & Kosinus-Ähnlichkeit \\
        \hline

        Anzahl Ergebnisse (Frontend)
            & 1 bis 5, Standard: 1 \\
        \hline
    \end{tabularx}
\end{table}

% ============================================================
% Prompt-Liste für YOLO-World
% ============================================================
Verwendete Prompts für die Merkmalsdetektion mit YOLO-World
\begin{table}[htbp]
    \centering
    \caption{Verwendete Prompts für die Merkmalsdetektion mit YOLO-World}
    \label{tab:promptliste}

    \begin{tabularx}{\textwidth}{|l|X|}
        \hline
        \textbf{Gruppe} &
        \textbf{Prompts} \\
        \hline

        Kopf und Gesicht (10)
            & person, human face, profile portrait, head, beard, hair, eye, nose, mouth, ear \\
        \hline

        Kopfschmuck (5)
            & crown, helmet, laurel wreath, wreath, diadem \\
        \hline

        Tiere (10)
            & bird, eagle, owl, horse, lion, bull, dog, wolf, snake, hedgehog \\
        \hline

        Waffen und Attribute (9)
            & shield, sword, spear, bow, arrow, torch, staff, scepter, trident \\
        \hline

        Pflanzen (5)
            & leaf, branch, tree, flower, vine \\
        \hline

        Architektur (5)
            & temple, building, column, tower, arch \\
        \hline

        Himmelskörper und Zeichen (5)
            & star, sun, moon, cross, circle \\
        \hline

        Schrift (4)
            & text, letter, symbol, inscription \\
        \hline

        Sonstiges (4)
            & coin, statue, mask, face \\
        \hline
    \end{tabularx}
\end{table}